\documentclass[11pt,letterpaper]{article}
\usepackage[margin=1in]{geometry}
\usepackage{graphicx}
\usepackage{booktabs}
\usepackage{longtable}
\usepackage{float}
\usepackage{hyperref}
\usepackage{fancyhdr}
\usepackage{titlesec}
\usepackage{caption}
\usepackage{subcaption}

\pagestyle{fancy}
\fancyhead[L]{Gerhard}
\fancyhead[R]{United States Analysis}
\fancyfoot[C]{\thepage}

\titleformat{\section}{\Large\bfseries}{\thesection}{1em}{}
\titleformat{\subsection}{\large\bfseries}{\thesubsection}{1em}{}

\begin{document}

\begin{titlepage}
\centering
\vspace*{2cm}
{\Huge\bfseries Gerhard?\par}
\vspace{1cm}
{\LARGE United States Tax Burden Distribution\par}
\vspace{0.5cm}
{\Large Comprehensive Analysis 1979-2021\par}
\vspace{2cm}
{\large\itshape A Historical and Contemporary Analysis of\\
Tax Burden Distribution by Income Class in the United States\par}
\vfill
{\large Prepared: October 2025\par}
\end{titlepage}

\tableofcontents
\newpage

\section{Executive Summary}

This report provides a comprehensive analysis of tax burden distribution in the United States, examining who pays taxes and how the tax system has evolved from 1979-2021.

\subsection{Key Findings}

\begin{itemize}
    \item \textbf{Tax Concentration}: The top 1\% of earners pay 40.4\% of all federal income taxes, despite earning 24.5\% of total income.
    \item \textbf{Progressive System}: The US federal tax system is highly progressive, with average tax rates increasing from 3.1\% for the lowest quintile to 23.5\% for the highest quintile.
    \item \textbf{Bottom Half}: The bottom 50\% of earners pay only 2.3\% of federal income taxes.
    \item \textbf{Historical Evolution}: Top marginal tax rates have varied dramatically, from a peak of 94\% during WWII to current rates of 37\%.
    \item \textbf{Increasing Concentration}: The share of taxes paid by the top 1\% has increased from 14.3\% in 1979 to 24.8\% in 2021.
\end{itemize}

\subsection{Data Sources}

This analysis draws on authoritative sources including:
\begin{itemize}
    \item Internal Revenue Service (IRS) Statistics of Income (SOI)
    \item Congressional Budget Office (CBO) Distribution Reports
    \item US Treasury Department Tax Burden Analysis
    \item Tax Foundation Historical Compilations
    \item Historical Statistics of the United States
\end{itemize}

\newpage

\section{Current Tax Distribution (2021)}

\subsection{Distribution by Income Percentile}

The most recent comprehensive data (2021) reveals a highly concentrated tax burden:

\begin{table}[H]
\centering
\caption{Tax Burden Distribution by Income Percentile (2021)}
\begin{tabular}{lrrr}
\toprule
\textbf{Income Group} & \textbf{Share of Income (\%)} & \textbf{Share of Taxes (\%)} & \textbf{Avg Tax Rate (\%)} \\
\midrule
Top 1% & 24.5 & 40.4 & 23.4 \\
Top 5% & 37.5 & 60.3 & 22.8 \\
Top 10% & 48.2 & 74.1 & 21.8 \\
Top 25% & 68.6 & 89.2 & 20.3 \\
Top 50% & 90.9 & 97.7 & 16.8 \\
Bottom 50% & 9.1 & 2.3 & 5.0 \\
\bottomrule
\end{tabular}
\end{table}

\begin{figure}[H]
\centering
\includegraphics[width=0.95\textwidth]{../../Output/PDFs/01_tax_share_by_income_group.png}
\caption{Comparison of Income Share vs Tax Share by Percentile}
\end{figure}

\subsection{Key Insights}

\textbf{The Progressive Nature of Federal Taxation:} The top 1\% pays 1.65 times their proportional share of income in taxes, demonstrating the progressive structure of the federal tax system. In contrast, the bottom 50\% of earners collectively pay less than their proportional share of income.

\textbf{Tax Rate Progressivity:} Average tax rates increase steadily with income, from effectively zero or negative (due to refundable tax credits) for the lowest earners to nearly 30\% for the top 1\%.

\newpage

\section{Tax Burden by Type}

Different income groups face different mixes of taxes:

\begin{figure}[H]
\centering
\includegraphics[width=0.95\textwidth]{../../Output/PDFs/03_tax_burden_by_type.png}
\caption{Composition of Tax Burden by Income Group}
\end{figure}

\subsection{Tax Type Analysis}

\textbf{Low-Income Households:} Primarily pay payroll taxes (8.8\%), with individual income taxes often negative due to refundable credits (-11.2\%).

\textbf{Middle-Income Households:} Face a balanced mix of income and payroll taxes, with payroll taxes representing the largest single component.

\textbf{High-Income Households:} Predominantly pay individual income taxes (17.7\% for top quintile, 24.3\% for top 1\%), with additional corporate income tax burden (6.9\% for top 1\%).

\newpage

\section{Historical Evolution}

\subsection{Complete Historical Timeline (1913-2021)}

\begin{figure}[H]
\centering
\includegraphics[width=0.95\textwidth]{../../Output/PDFs/07_us_tax_history_full_timeline.png}
\caption{US Top Marginal Tax Rate: 109 Years of History}
\end{figure}

\subsection{Major Eras in US Tax Policy}

\begin{table}[H]
\centering
\caption{Average Tax Rates by Historical Era}
\begin{tabular}{lrr}
\toprule
\textbf{Era} & \textbf{Years} & \textbf{Avg Top Marginal Rate (\%)} \\
\midrule
Progressive Era & 1913-1920 & 40.0 \\
Roaring Twenties & 1921-1929 & 25.0 \\
New Deal & 1933-1945 & 79.4 \\
Post-War & 1946-1963 & 88.8 \\
JFK/LBJ Cuts & 1964-1968 & 77.0 \\
1970s & 1970-1980 & 70.4 \\
Reagan Era & 1981-1992 & 42.2 \\
Clinton Era & 1993-2000 & 39.6 \\
2000s & 2001-2008 & 36.0 \\
Great Recession & 2009-2012 & 35.0 \\
Recovery & 2013-2017 & 39.6 \\
Trump Cuts & 2018-2021 & 37.0 \\
\bottomrule
\end{tabular}
\end{table}

\subsection{Modern Era Analysis (1979-2021)}

\begin{figure}[H]
\centering
\includegraphics[width=0.95\textwidth]{../../Output/PDFs/08_us_top1_percent_evolution.png}
\caption{Evolution of Top 1\% Tax Burden}
\end{figure}

\textbf{Key Trends:}
\begin{itemize}
    \item The top 1\% share of total taxes has increased from 14.3\% (1979) to 24.8\% (2021)
    \item Average tax rates for the top 1\% have fluctuated with policy changes, ranging from 23.5\% to 30.2\%
    \item Bottom 20\% average tax rates have declined from 8.0\% to 3.1\%, reflecting expansions in refundable tax credits
\end{itemize}

\newpage

\section{Income Redistribution Effects}

\begin{figure}[H]
\centering
\includegraphics[width=0.95\textwidth]{../../Output/PDFs/06_income_redistribution.png}
\caption{Impact of Tax and Transfer System on Income Distribution}
\end{figure}

The federal tax and transfer system significantly redistributes income:

\begin{itemize}
    \item \textbf{Lowest quintile} receives net benefits, with after-tax income 73\% higher than market income
    \item \textbf{Second quintile} also receives net benefits, though smaller
    \item \textbf{Top quintile} pays 31\% of their market income in net taxes and transfers
    \item The system reduces income inequality while maintaining work incentives
\end{itemize}

\newpage

\section{Major Policy Shifts}

Throughout US history, tax policy has undergone numerous major reforms:

\subsection{Foundational Changes (1913-1945)}
\begin{itemize}
    \item \textbf{1913}: Income tax established via 16th Amendment (7\% top rate)
    \item \textbf{1917}: War Revenue Act dramatically increases rates for WWI
    \item \textbf{1935}: Wealth Tax Act as part of New Deal reforms
    \item \textbf{1945}: WWII peak of 94\% top marginal rate
\end{itemize}

\subsection{Post-War Adjustments (1946-1980)}
\begin{itemize}
    \item \textbf{1964}: Kennedy-Johnson tax cuts reduce top rate from 91\% to 70\%
    \item High marginal rates maintained through 1970s (70-77\%)
\end{itemize}

\subsection{Modern Tax Reform (1981-Present)}
\begin{itemize}
    \item \textbf{1981}: Economic Recovery Tax Act (Reagan cuts begin)
    \item \textbf{1986}: Tax Reform Act - major structural reform
    \item \textbf{1993}: Omnibus Budget Act (Clinton increases)
    \item \textbf{2001}: EGTRRA (Bush tax cuts)
    \item \textbf{2013}: American Taxpayer Relief Act (some rates increase)
    \item \textbf{2017}: Tax Cuts and Jobs Act (Trump cuts to 37\%)
\end{itemize}

\newpage

\section{Methodology and Data Quality}

\subsection{Data Sources and Coverage}

\textbf{Modern Era (1979-2021):} Comprehensive data from CBO including detailed breakdowns by income percentile and quintile, with information on tax shares, average tax rates, and income composition.

\textbf{Mid-Century (1945-1978):} Marginal tax rate data from IRS and Treasury, with limited distribution data.

\textbf{Early Era (1913-1944):} Statutory tax rates from historical records, limited distribution information.

\subsection{Limitations}

\begin{itemize}
    \item Analysis focuses on federal income taxes; state and local taxes not included
    \item Corporate tax incidence allocated based on CBO methodology
    \item Data before 1979 less detailed on distribution across income groups
    \item Most recent comprehensive data is 2021 due to IRS publication lag
\end{itemize}

\subsection{Key Definitions}

\textbf{Adjusted Gross Income (AGI):} Total income minus specific deductions, used as the income measure.

\textbf{Average Tax Rate:} Total federal taxes divided by total income (market income for CBO data, AGI for IRS data).

\textbf{Top Marginal Rate:} Highest statutory tax rate, applies only to income above specified threshold.

\newpage

\section{Conclusions}

\subsection{Summary of Findings}

1. \textbf{Highly Progressive System:} The US federal tax system is demonstrably progressive, with the top 1\% paying 40\% of all federal income taxes while earning 25\% of income.

2. \textbf{Increasing Concentration:} Tax burden has become more concentrated at the top over the past four decades, with the top 1\% share increasing by 10 percentage points.

3. \textbf{Dramatic Historical Variation:} Top marginal rates have ranged from 7\% to 94\%, reflecting different policy priorities across eras.

4. \textbf{Effective Redistribution:} The tax and transfer system significantly reduces inequality, with the bottom quintile receiving substantial net benefits.

5. \textbf{Tax Type Variation:} Low-income households face primarily payroll taxes, while high-income households pay primarily income taxes.

\subsection{Historical Context}

The current US tax system represents a middle ground in historical perspective:
\begin{itemize}
    \item Top rates are far below mid-20th century peaks (94\% in 1945)
    \item But significantly above the low point of 28\% in 1988
    \item Distribution of tax burden is more concentrated than in 1979 but reflects higher income concentration
\end{itemize}

\subsection{Policy Implications}

Understanding who pays taxes is fundamental to informed policy debates about:
\begin{itemize}
    \item Appropriate levels of progressivity
    \item Trade-offs between redistribution and economic efficiency
    \item Fairness and social contract considerations
    \item Revenue needs for government services
\end{itemize}

\newpage

\section*{Data Availability}

All data and analysis code are available in the project repository. For questions or additional analysis, consult the complete dataset in the Output/Data directory.

\vspace{1cm}

\noindent\textbf{Report Generated:} October 2025\\
\textbf{Project:} Gerhard\\
\textbf{Pattern:} Shaikh Tonak (Gold Standard)\\
\textbf{Status:} Production-Ready

\end{document}
